% ═══════════════════════════════════════════════════════════════════════════
% QR-CODE: Lernpfad Metalloxide (Chemie Klasse 8)
% ═══════════════════════════════════════════════════════════════════════════

\documentclass[a4paper,11pt]{article}

\usepackage[utf8]{inputenc}
\usepackage[T1]{fontenc}
\usepackage{helvet}
\renewcommand{\familydefault}{\sfdefault}

\usepackage[margin=20mm]{geometry}
\usepackage{tikz}
\usepackage{xcolor}
\usepackage{qrcode}

% ═══════════════════════════════════════════════════════════════════════════
% FARBEN
% ═══════════════════════════════════════════════════════════════════════════

\definecolor{chemie}{HTML}{2a7a4b}
\colorlet{fachfarbe}{chemie}

% ═══════════════════════════════════════════════════════════════════════════
% KONFIGURATION
% ═══════════════════════════════════════════════════════════════════════════

\newcommand{\lptitel}{Reaktion von Metallen mit Sauerstoff}
\newcommand{\lpurl}{https://devbox42.github.io/sciencesim/chemie/kl08/01-metalloxide/LP-01-metalloxide.html}
\newcommand{\lpurlkurz}{devbox42.github.io/.../LP-01-metalloxide.html}

% ═══════════════════════════════════════════════════════════════════════════
% DOKUMENT
% ═══════════════════════════════════════════════════════════════════════════

\pagestyle{empty}

\begin{document}

\begin{center}

\vspace*{10mm}
{\fontsize{28}{34}\selectfont\bfseries\color{fachfarbe}Lernpfad: \lptitel}

\vspace{15mm}

\fcolorbox{fachfarbe}{white}{\qrcode[height=100mm]{\lpurl}}

\vspace{12mm}

{\large\color{gray}\lpurlkurz}

\vspace{8mm}

{\Large\color{gray}Scanne den QR-Code mit deinem Gerät}

\vspace{10mm}

\fcolorbox{fachfarbe}{fachfarbe!10}{%
    \parbox{0.7\textwidth}{%
        \centering
        \vspace{3mm}
        Arbeite die 8 Schritte durch und fülle dabei\\[2mm]
        das Arbeitsblatt aus.
        \vspace{3mm}
    }%
}

\end{center}

\end{document}
